\section{Preliminaries}\label{section2}

\subsection{Star-product quantization of a graded Poisson algebra}
Let $A$ be a commutative ${\mathbb C}$-algebra with a $\Z_{\geq 0}$ grading given by
\[
A = \bigoplus_{d \geq 0} A_d.
\]
Let $\{~,~\}$ be a Poisson bracket of degree $-2$ on $A$; namely,
\[
\{~,~\} \colon \ A_n \otimes A_m \to A_{n+m-2}.
\]
A basic example is the symmetric algebra $A=SV$ of a finite dimensional symplectic vector space $V$
with the Poisson bracket defined by the symplectic form.

Note that this setting includes the case of graded Poisson algebras with Poisson bracket of degree $-1$ (such as the symmetric algebra of a Lie algebra, the algebra of regular functions on the nilpotent cone, etc.), by multiplying the degrees by $2$.\footnote{In fact, this theory can be developed for Poisson brackets of any negative degree, but for most applications it is sufficient to consider degree $-2$, so we restrict our attention to this case.}

\begin{Definition}
A \emph{star-product} on $A$ quantizing $\lbrace\,,\,\rbrace$ is an associative multiplication operation
\[
*\colon \ A \otimes A \to A
\]
such that for $a \in A_n$ and $b \in A_m$,
\[
a*b = \sum_{k = 0}^{\lfloor \frac{n+m}{2}\rfloor} C_k(a,b),
\]
where $C_k \colon A_n \otimes A_m \to A_{n+m-2k}$ are bilinear maps such that $C_0(a,b) = a b$ and
\begin{gather*}
C_1(a,b) - C_1(b,a) = \{a,b\}.
\end{gather*}
\end{Definition}

\begin{Definition} A star-product $*$ on $A$ is \emph{short} if
for any $m,n\in {\mathbb Z}_{\ge 0}$ and any $a \in A_n$, $b \in A_m$ one has $C_k(a,b) = 0$ for all $k > \min(n,m)$. In other words, $a * b$ has no terms in $A_d$ for $d < |n - m|$.
\end{Definition}

From now on assume that $A_0={\mathbb C}$ and $\dim A_i<\infty$ for all $i$. We define the following (in general, non-symmetric) inner product on $A$ induced by the star-product. Let $a \in A_n$ and $b \in A_m$. Then
\[
\langle a, b \rangle := \text{CT}(a * b),
\]
where $\text{CT}\colon A \to A_0$ is the constant term map, taking the term of degree zero. Note that if $*$ is short, then $\langle a, b \rangle = 0$ if $n \neq m$, i.e., the decomposition $A=\oplus_{d\ge 0}A_d$ is orthogonal.

\begin{Definition} A short star-product $*$ on $A$ is said to be {\it nondegenerate} if the corresponding inner product $\langle \,,\,\rangle$ is nondegenerate in each degree~$i$.
\end{Definition}

\subsection{Construction of star-products from a quantum algebra}\label{constru} Let $A$ be a ${\mathbb Z}_{\ge 0}$-graded algebra. Fix a filtered associative algebra
\[
{{\mathbf A}} = \bigcup_{d \geq 0} F_d{{\mathbf A}}
\]
such that its associated graded algebra is identified with the graded algebra $A$; namely,
\[
\operatorname{gr}{{\mathbf A}} = \bigoplus_{d \geq 0} F_{d}{{\mathbf A}}/F_{d-1}{{\mathbf A}}
\]
and $F_{d}{{\mathbf A}}/F_{d-1}{{\mathbf A}} \cong A_d$ for $d \geq 0$ (compatibly with multiplication) with $F_{-1}{{\mathbf A}} := 0$.
Assume that we have a filtration preserving ${\mathbb Z}/2$-action $s\colon {{\mathbf A}}\to {{\mathbf A}}$ such that the corresponding associated graded map $s\colon A\to A$ is given by $s=(-1)^d$, where $d\colon A\to A$ is the degree operator. In this case, it is easy to see that $[F_n{{\mathbf A}},F_m{{\mathbf A}}]\subset F_{n+m-2}{{\mathbf A}}$, so we have the leading coefficient map $\lbrace\,,\,\rbrace\colon A_n\otimes A_m\to A_{n+m-2}$. It is easy to check that $\lbrace\,,\,\rbrace$ is a Poisson bracket on $A$, and ${{\mathbf A}}$ is a~${\mathbb Z}/2$-equivariant filtered quantization of $\lbrace\,,\,\rbrace$.\footnote{Since all quantizations we will consider will be ${\mathbb Z}/2$-equivariant, we will not mention it explicitly from now on.}

\begin{Definition}A \emph{quantization map} is a ${\mathbb Z}/2$-equivariant isomorphism of vector spaces
\[
\phi \colon \ A \to {{\mathbf A}}
\]
such that $\phi(A_d) \subseteq F_d{{\mathbf A}}$ for $d\ge 0$ and $\operatorname{gr} \phi \colon A \to A$ is the identity map, where
\[
\operatorname{gr} \phi = \bigoplus_{d\geq 0} \operatorname{gr} \phi_d, \qquad \operatorname{gr} \phi_d \colon \ A_d \to F_d{{\mathbf A}}/F_{d-1}{{\mathbf A}}=A_d.
\]
\end{Definition}
Note that $\phi$ need not (and usually cannot) be a homomorphism of algebras, since the algebra~$A$ is commutative but ${{\mathbf A}}$ is not in general.

Given a quantization map $\phi\colon A \to {{\mathbf A}}$, we can define the star-product on $A$ by
\[
a * b := \phi^{-1}(\phi(a)\phi(b)).
\]
Conversely, given a star-product on $A$, we can set ${{\mathbf A}}:=(A,*)$ with the filtration induced by the grading, and $\phi:={\rm Id}$.

\begin{Proposition} This defines a pair of mutually inverse bijections between star-products on~$A$ and ${\mathbb Z}/2$-equivariant quantizations of $A$ equipped with a quantization map.
\end{Proposition}

\begin{proof} Straightforward.
\end{proof}

\subsection{Even star-products}

\begin{Definition} A star-product $*$ on $A$ is said to be {\it even} if $C_k(a,b)=(-1)^k C_k(b,a)$ for all $k\ge 0$.
\end{Definition}

In particular, for even star-products $C_1$ is skew-symmetric, so since $C_1(a,b)-C_1(b,a)=\lbrace a,b\rbrace$, we have $C_1(a,b)=\frac{1}{2}\lbrace a,b\rbrace$.

Let $*$ be an even star-product on $A$ and ${{\mathbf A}}=(A,*)$ be the corresponding quantum algebra.
Then the map $\sigma={\rm i}^d$ defines an antiautomorphism ${{\mathbf A}}\to {{\mathbf A}}$ such that $\sigma^2=s=(-1)^d$. A ${\mathbb Z}/2$-invariant quantization ${\mathbf A}$ of $A$ equipped with a~filtration-preserving antiautomorphism $\sigma$ such that $\sigma^2=s$ will be called {\it even}. Thus an even star-product gives rise to an even quantization. Conversely, an even quantization ${{\mathbf A}}$ of $A$ gives rise to an even star-product on~$A$.

\begin{Proposition} This defines a pair of mutually inverse bijections between even star-products on $A$ and ${\mathbb Z}/2$-equivariant quantizations ${{\mathbf A}}$ of $A$ equipped with a filtered antiautomorphism $\sigma\colon {{\mathbf A}}\to {{\mathbf A}}$ such that $\operatorname{gr}\sigma={\rm i}^d$ and $\sigma^2=s$ with a $\sigma$-invariant $($i.e., ${\mathbb Z}/4$-invariant$)$ quantization map.
\end{Proposition}

\begin{proof} Straightforward.
\end{proof}
